\section{Lidar application details}
\label{app:lidar-computation}

\subsection{Computation}

For the model in \cref{sec:lidar}, computation subtracts the observed
response mean and divides by the response range, applying the same
change of scale to the Gaussian model. The global-threshold tolerance
is $10^{-12}$ in these transformed coefficient units. An expectation
propagation approximation at $129$ derivative locations supplies the
initial mean and dense proposal metric. The continuum threshold
still determines every target evaluation, so these locations affect
only initialization and the metric.

Four Stan chains each use $400$ warm-up iterations and $2{,}000$
retained draws, without thinning. During warm-up, the step size adapts
to target acceptance $0.9$ while the supplied metric is held fixed.
The maximum tree depth is $10$ during warm-up and $12$ during
retained sampling, which otherwise uses the resulting tuning.

\paragraph{MCMC diagnostics.}
We monitor the sampling coordinates and reconstructed coefficients,
the latent response and concentration at the observed ranges, their
adjacent midpoints, and five equally spaced locations, together with
the total response increase, prior quadratic form, and log density.
Across these $1{,}146$ scalar quantities, the maximum rank-normalized
split $\widehat R$ is $1.0034$, and the minimum bulk and tail effective
sample sizes are $2{,}776$ and $2{,}479$, respectively. None of the
$8{,}000$ retained iterations has a divergence or reaches the
tree-depth limit, and the minimum energy Bayesian fraction of missing
information (E-BFMI) across chains is $0.9637$. Three divergences
occurred during warm-up, whose draws were discarded. The retained-draw
diagnostics reveal no warning about mixing in the monitored quantities,
while the predictive checks below assess agreement with the observation
model.

\subsection{Predictive checks}

For each retained monotone posterior curve, define the standardized residuals $e_i(\boldsymbol{a})=(y_i-f_{\boldsymbol{a}}(x(r_i)))/\sigma_i$.
We generate one independent replicate dataset from
\[
 Y_i^{\mathrm{rep}}\mid \boldsymbol{a}\overset{\mathrm{ind}}\sim N(f_{\boldsymbol{a}}(x(r_i)),\sigma_i^2),
 \qquad e_i^{\mathrm{rep}}(\boldsymbol{a})=\frac{Y_i^{\mathrm{rep}}-f_{\boldsymbol{a}}(x(r_i))}{\sigma_i}.
\]
We compute the mean squared residual $n^{-1}\sum_i e_i(\boldsymbol{a})^2$, maximum $\max_i|e_i(\boldsymbol{a})|$, and absolute sample correlation between the paired vectors $(e_1,\ldots,e_{n-1})$ and $(e_2,\ldots,e_n)$, ordered by recorded distance.
Here adjacency refers to consecutive released ranges, whose spacings vary. The two remaining discrepancies are the signed mean residuals over $r_i\leq555$ meters and $r_i>555$ meters, each averaged over the observations in its own half of the range. Applying the same definitions to $e_i^{\mathrm{rep}}$ gives the replicated discrepancies in \cref{tab:lidarcriticism}.

\begin{table}[htbp]
\centering
\begin{tabular}{lrr}
\toprule
Constrained posterior discrepancy & Lower tail & Upper tail \\
\midrule
Mean squared standardized residual & 0.5931 & 0.4069 \\
Maximum absolute standardized residual & 0.3914 & 0.6086 \\
Absolute adjacent residual correlation & 0.5359 & 0.4641 \\
Mean standardized residual at $r\leq555$ m & 0.3526 & 0.6474 \\
Mean standardized residual at $r>555$ m & 0.4235 & 0.5765 \\
\bottomrule
\end{tabular}
\caption[Residual checks for mercury lidar]{
Posterior predictive checks from $8{,}000$ globally monotone posterior curves, with one independent Gaussian replicate dataset per curve.
Lower and upper tails give the fractions of replicated discrepancies at or below and at or above the observed discrepancies for the same posterior curve, with the two range summaries split at the midpoint of $555$ meters.
The summaries condition on the fitted covariance and noise parameters and reuse the fitting data, so they are neither calibrated test probabilities nor out-of-sample predictive scores.}
\label{tab:lidarcriticism}
\end{table}

The upper-tail fractions in \cref{tab:lidarcriticism} range from
$0.4069$ to $0.6474$, with no pronounced disagreement for these
summaries under the specified observation model.

\FloatBarrier
